
\section{Đề xuất kiến trúc phân mảnh trạng thái}

Máy giặt là một hệ thống động lực học có tính chu kỳ, trong đó các pha vận hành mang đặc tính vật lý (âm thanh và biên độ rung) thay đổi rõ rệt. Việc sử dụng một mô hình mạng nơ-ron tự mã hóa duy nhất để bao quát toàn bộ chu kỳ sẽ làm tăng độ phức tạp của không gian phân phối dữ liệu, dẫn đến tỷ lệ cảnh báo giả (False Positive) cao tại các thời điểm thiết bị chuyển pha. Các nghiên cứu trước đây thường chỉ tập trung giám sát trên một trạng thái tĩnh nhất định, bỏ qua vấn đề này \cite{ref_1_lord}. Để giải quyết bài toán này, đồ án đề xuất kiến trúc mạng nơ-ron tự mã hóa Ba trạng thái.

\subsection{Cơ sở phân tách chu trình hoạt động}

Để hệ thống có thể tự xác định trạng thái vận hành cục bộ mà không cần can thiệp vào mạch điều khiển trung tâm của máy giặt, một thuật toán phân luồng trạng thái dựa trên dữ liệu chấn động đã được thiết kế.

Thông qua phân tích tương quan thực nghiệm trên tập dữ liệu thu thập, nghiên cứu xác định được \textbf{phương sai trục Z} (đại diện cho cường độ chấn động thẳng đứng) của cảm biến ADXL345 là chỉ báo ổn định để phân loại trạng thái. Logic định tuyến được xây dựng bằng cách so sánh với các ngưỡng phân hoạch tĩnh, chia chu kỳ thành ba giai đoạn chính: \textbf{GENTLE} (Giặt/Ngâm), \textbf{STRONG} (Xả/Đảo lồng) và \textbf{SPIN} (Vắt tốc độ cao).

Các ngưỡng này được trích xuất từ phân tích phân phối thống kê của tập dữ liệu huấn luyện (trình bày chi tiết tại Chương 3). Cơ chế định tuyến này cho phép tác vụ \texttt{ai\_task} tự động trỏ con trỏ bộ thông dịch (\texttt{Interpreter}) tới vùng nhớ chứa trọng số mô hình tương ứng. Do cả ba mô hình đã được nạp sẵn trên PSRAM, quá trình chuyển đổi diễn ra tức thời, đáp ứng yêu cầu thời gian thực của hệ thống.

\subsection{Thiết kế mạng tự mã hóa Dense chuyên biệt}

Mỗi mô hình con được xây dựng theo kiến trúc mạng nơ-ron tự mã hóa Dense với cấu trúc đối xứng, giúp đơn giản hóa tính toán và cho phép tái sử dụng bộ đệm hiệu quả trong môi trường nhúng:

\begin{figure}[H]
\centering
\resizebox{0.98\textwidth}{!}{%
\begin{tikzpicture}[
    layer/.style={draw, rounded corners=2pt, align=center, minimum width=1.55cm, minimum height=1.0cm, font=\small},
    arrow/.style={->, thick}
]
\node[layer, fill=blue!8] (in) at (0,0) {Đầu vào\\19};
\node[layer, fill=green!10] (e1) at (2.0,0) {Dense\\128};
\node[layer, fill=green!10] (e2) at (4.0,0) {Dense\\64};
\node[layer, fill=orange!18] (b) at (6.0,0) {Cổ chai\\32};
\node[layer, fill=green!10] (d1) at (8.0,0) {Dense\\64};
\node[layer, fill=green!10] (d2) at (10.0,0) {Dense\\128};
\node[layer, fill=blue!8] (out) at (12.0,0) {Đầu ra\\19};
\draw[arrow] (in) -- (e1);
\draw[arrow] (e1) -- (e2);
\draw[arrow] (e2) -- (b);
\draw[arrow] (b) -- (d1);
\draw[arrow] (d1) -- (d2);
\draw[arrow] (d2) -- (out);
\node[font=\small] at (3.0,-0.9) {Khối mã hóa};
\node[font=\small] at (9.0,-0.9) {Khối giải mã};
\end{tikzpicture}
}
\caption{Kiến trúc mạng tự mã hóa Dense đối xứng dùng cho mỗi pha vận hành.}
\label{fig:symmetric_dense_ae_arch}
\end{figure}

\begin{itemize}
    \item \textbf{Lớp đầu vào:} Tiếp nhận vector đặc trưng 19 chiều (13 đặc trưng Mel-filterbank và 6 đặc trưng rung động).
    \item \textbf{Khối mã hóa:} Thực hiện nén dữ liệu qua các lớp ẩn ($19 \rightarrow 128 \rightarrow 64 \rightarrow 32$). Lớp ẩn đầu tiên được thiết kế theo hướng biểu diễn vượt quá nhằm tăng cường khả năng trích xuất các tương quan phi tuyến từ luồng tín hiệu đa phương thức \cite{ref_2_overcomplete}.
    \item \textbf{Khối giải mã:} Có cấu trúc phản chiếu đối xứng ($32 \rightarrow 64 \rightarrow 128 \rightarrow 19$), làm nhiệm vụ tái tạo lại vector đầu vào ban đầu.
\end{itemize}

\subsection{Kỹ thuật cân bằng đặc trưng với hàm mất mát MAE có trọng số}

Trong kỹ thuật kết hợp dữ liệu đa phương thức, sự chênh lệch về số lượng đặc trưng (13 chiều âm thanh so với 6 chiều rung động) có thể khiến thuật toán học máy ưu tiên giảm thiểu sai số của luồng âm học và bỏ qua các biến động cơ học nhỏ. Để khắc phục, đồ án áp dụng kỹ thuật hàm mất mát MAE có trọng số trong giai đoạn huấn luyện. Hàm mất mát được định nghĩa như sau:
\begin{equation}
    L_{WMAE} = \frac{1}{N} \sum_{i=1}^{N} w_i \cdot |x_i - \hat{x}_i|
\end{equation}
Trong đó, hệ số $w_i$ được đặt bằng $1$ đối với các đặc trưng âm thanh và bằng $5$ đối với các đặc trưng rung động. Giá trị 5 được lựa chọn theo hướng thực nghiệm có kiểm soát: nó bù lại bất lợi về số chiều của nhánh rung động (6 chiều so với 13 chiều âm thanh) nhưng vẫn không làm triệt tiêu hoàn toàn đóng góp của phổ âm thanh. Bằng cách tăng trọng số phạt trên sai số tái tạo của rung động, mô hình được định hướng biểu diễn chính xác hơn các đặc trưng cơ học, từ đó cải thiện độ nhạy với các sự cố bất thường vật lý. Hiệu quả của lựa chọn này được kiểm chứng trong nghiên cứu loại bỏ thành phần ở Chương 4, khi biến thể không dùng MAE có trọng số cho kết quả thấp hơn rõ rệt so với mô hình đề xuất.

\subsection{Đánh giá tương thích hàm kích hoạt trong môi trường lượng tử hóa}

Để triển khai hiệu quả trên vi điều khiển Seeed Studio XIAO ESP32-S3, phương pháp lượng tử hóa nhận thức huấn luyện được tích hợp. Kỹ thuật này ánh xạ toàn bộ trọng số mạng nơ-ron từ định dạng số thực Float32 sang định dạng số nguyên INT8.

Đặc biệt, tại lớp thắt cổ chai, đồ án sử dụng hàm kích hoạt Sigmoid thay cho hàm ReLU truyền thống. Quyết định thiết kế này dựa trên đặc tính hữu hạn của Sigmoid (giới hạn đầu ra trong khoảng $[0, 1]$). Khoảng giá trị này giúp quá trình lượng tử hóa ánh xạ đồng đều hơn vào 256 mức rời rạc của chuẩn INT8, qua đó giảm sai số mất mát thông tin so với hàm ReLU có dải giá trị dương không giới hạn \cite{ref_1_14}.

% ==============================================================================
% DANH MỤC TÀI LIỆU THAM KHẢO TRÍCH DẪN (CHI TIẾT CHO MỤC 2.4)
% ==============================================================================
% \bibitem{ref_1_lord} M. Lord, "TinyML Anomaly Detection," Master's thesis, CSUN, 2021.
% \bibitem{ref_2_overcomplete} Y. Bengio, A. Courville, and P. Vincent, "Representation Learning: A Review and New Perspectives," IEEE Trans. Pattern Anal. Mach. Intell., vol. 35, no. 8, pp. 1798-1828, 2013.
% \bibitem{ref_1_14} B. Jacob et al., "Quantization and Training of Neural Networks for Efficient Integer-Arithmetic-Only Inference," CVPR, 2018.
